% Mathématiques 5e — cours V3
% Statistiques et représentations de données

\section{Objectifs}

\begin{itemize}
\item Recueillir et organiser des données.
\item Calculer des effectifs.
\item Calculer des fréquences sous forme fractionnaire, décimale ou en pourcentage.
\item Lire et interpréter tableaux, diagrammes et graphiques.
\item Représenter des données sous forme de tableau, diagramme ou graphique cartésien.
\item Choisir une représentation adaptée à l'information à mettre en avant.
\item Calculer et interpréter la moyenne simple d'une série.
\item Utiliser un tableur-grapheur pour organiser et représenter des données.
\end{itemize}

\section{1. De la question aux données}

Une étude statistique commence par une question.

On peut par exemple chercher :
\[
\text{« Combien de temps les élèves mettent-ils pour venir au collège ? »}
\]

Il faut ensuite recueillir les réponses de manière organisée.

\begin{methode}
Avant de calculer quoi que ce soit :
\begin{enumerate}
\item identifier la population étudiée ;
\item préciser la donnée recueillie ;
\item choisir une unité si la donnée est mesurée ;
\item consigner toutes les observations.
\end{enumerate}
\end{methode}

\section{2. Effectif}

L'effectif d'une valeur ou d'une catégorie est le nombre d'individus correspondants.

\begin{exemple}
Sur $30$ élèves :
\[
12
\]
viennent en bus.

L'effectif de la catégorie « bus » est :
\[
\boxed{12}.
\]
\end{exemple}

L'effectif total est la somme des effectifs de toutes les catégories.

\section{3. Fréquence}

La fréquence d'une valeur ou d'une catégorie est :
\[
\boxed{
f=\dfrac{\text{effectif de la catégorie}}{\text{effectif total}}
}.
\]

Dans l'exemple précédent :
\[
f=\dfrac{12}{30}.
\]

On simplifie :
\[
\dfrac{12}{30}=\dfrac25.
\]

En écriture décimale :
\[
f=0{,}4.
\]

En pourcentage :
\[
f=40\%.
\]

\section{4. Vérifier les fréquences}

La somme des fréquences de toutes les catégories doit être égale à :
\[
1
\]
ou :
\[
100\%.
\]

Une petite différence peut apparaître si des pourcentages ont été arrondis.

\begin{attention}
Une fréquence ne peut pas être négative ni supérieure à $1$.
\end{attention}

\section{5. Tableau de données}

Un tableau permet d'organiser les valeurs et leurs effectifs.

\begin{tabular}{c|c|c}
Mode de transport & Effectif & Fréquence \\
Bus & $12$ & $40\%$ \\
Vélo & $8$ & $26{,}7\%$ \\
Marche & $6$ & $20\%$ \\
Voiture & $4$ & $13{,}3\%$ \\
\end{tabular}

Le tableau doit avoir :
\begin{itemize}
\item des intitulés clairs ;
\item les unités utiles ;
\item un total cohérent.
\end{itemize}

\section{6. Plusieurs représentations}

Une même série peut être présentée de plusieurs façons.

\coursefigure{/images/mathematiques/5eme/stats-representations.svg}{Tableau, diagramme en barres et diagramme circulaire représentant une même série.}{Le choix d'une représentation dépend de la question posée et du message que l'on souhaite mettre en évidence.}

Chaque représentation possède des avantages différents.

\section{7. Diagramme en barres}

Un diagramme en barres convient bien pour comparer des catégories.

La hauteur de chaque barre est proportionnelle à la valeur représentée.

Pour le lire, on vérifie :
\begin{itemize}
\item le titre ;
\item les catégories ;
\item l'axe des valeurs ;
\item l'échelle.
\end{itemize}

\section{8. Diagramme circulaire}

Un diagramme circulaire représente la répartition d'un total.

Le disque entier correspond à :
\[
100\%.
\]

Un secteur représentant :
\[
25\%
\]
occupe un quart du disque.

Un secteur représentant :
\[
50\%
\]
occupe la moitié.

\section{9. Graphique cartésien}

Un graphique cartésien permet souvent de représenter une évolution.

Par exemple :
\[
\text{température en fonction du temps}.
\]

L'axe horizontal peut représenter le temps.

L'axe vertical représente alors la température.

Il faut toujours lire les axes avant d'interpréter la courbe.

\section{10. Choisir une représentation}

On ne choisit pas un graphique uniquement parce qu'il est joli.

\begin{methode}
\begin{itemize}
\item Comparer des catégories : diagramme en barres.
\item Montrer une répartition : diagramme circulaire.
\item Montrer une évolution : graphique cartésien.
\item Donner des valeurs exactes : tableau.
\end{itemize}
\end{methode}

\section{11. Construire un diagramme en barres}

Supposons les effectifs :
\[
5,\quad8,\quad12,\quad3.
\]

On :
\begin{enumerate}
\item choisit une échelle ;
\item place les catégories ;
\item trace des barres de même largeur ;
\item donne une hauteur proportionnelle à chaque effectif ;
\item ajoute titre et légendes.
\end{enumerate}

\section{12. Construire un diagramme circulaire}

Pour une catégorie de fréquence :
\[
30\%,
\]
l'angle du secteur vaut :
\[
0{,}30\times360^\circ.
\]

Donc :
\[
108^\circ.
\]

En 5e, on peut exploiter ce lien pour construire ou vérifier une représentation.

\section{13. Lire une information}

Un graphique peut répondre à des questions immédiates :
\begin{itemize}
\item quelle catégorie est la plus fréquente ?
\item quelle valeur est maximale ?
\item entre quels instants la grandeur augmente-t-elle ?
\item quelles valeurs dépassent un seuil ?
\end{itemize}

La réponse doit être reliée à l'échelle et à l'unité.

\section{14. Lecture critique}

Un graphique peut accentuer visuellement une différence.

Un axe qui commence près des valeurs observées rend les écarts plus impressionnants.

Il faut donc contrôler :
\begin{itemize}
\item le zéro éventuel ;
\item les graduations ;
\item la régularité de l'échelle ;
\item les unités ;
\item la population étudiée.
\end{itemize}

\section{15. Moyenne simple}

La moyenne d'une série de $n$ valeurs est :
\[
\boxed{
\overline{x}
=
\dfrac{x_1+x_2+\cdots+x_n}{n}
}.
\]

\begin{exemple}
Pour les valeurs :
\[
8,\quad10,\quad12,\quad14,
\]
la somme vaut :
\[
44.
\]

Il y a :
\[
4
\]
valeurs.

La moyenne vaut :
\[
\dfrac{44}{4}=11.
\]

Donc :
\[
\boxed{\overline{x}=11}.
\]
\end{exemple}

\section{16. Interpréter la moyenne}

La moyenne peut être comprise comme une valeur de répartition équitable.

\coursefigure{/images/mathematiques/5eme/stats-moyenne.svg}{Colonnes de hauteurs différentes rééquilibrées en colonnes de même hauteur.}{La moyenne conserve le total tout en répartissant la somme de manière égale entre toutes les valeurs.}

Si un total de :
\[
44
\]
est réparti également entre :
\[
4
\]
valeurs, chacune reçoit :
\[
11.
\]

\section{17. Moyenne et unité}

La moyenne possède la même unité que les valeurs.

Si les valeurs sont des durées en minutes, la moyenne s'exprime en :
\[
\text{minutes}.
\]

Si elles sont des températures en degrés Celsius, la moyenne s'exprime en :
\[
^\circ\text{C}.
\]

\section{18. Moyenne non présente dans la série}

La moyenne n'est pas obligée d'être une valeur réellement observée.

Par exemple, pour :
\[
10,\quad11,
\]
la moyenne vaut :
\[
10{,}5.
\]

La valeur :
\[
10{,}5
\]
n'apparaît pourtant pas dans la série.

\section{19. Exemple développé}

Les durées de cinq trajets sont :
\[
12,\quad15,\quad9,\quad14,\quad10\ \text{min}.
\]

La somme vaut :
\[
12+15+9+14+10=60.
\]

La moyenne est :
\[
\dfrac{60}{5}=12.
\]

Donc :
\[
\boxed{12\ \text{min}}.
\]

\section{20. Comparer deux séries}

On peut comparer :
\begin{itemize}
\item leurs effectifs ;
\item leurs fréquences ;
\item leurs représentations ;
\item leurs moyennes.
\end{itemize}

Une moyenne identique ne signifie pas que les séries sont identiques.

Par exemple :
\[
8,\quad12
\]
et :
\[
10,\quad10
\]
ont toutes deux pour moyenne :
\[
10.
\]

\section{21. Utiliser un tableur}

Un tableur permet :
\begin{itemize}
\item de saisir des données ;
\item de compter ;
\item de calculer des fréquences ;
\item de calculer une moyenne ;
\item de produire des graphiques.
\end{itemize}

Le tableur n'interprète pas automatiquement la pertinence du résultat.

L'élève doit encore choisir les données et la représentation.

\section{22. Méthode générale}

\begin{methode}
Pour étudier une série :
\begin{enumerate}
\item organiser les données ;
\item calculer les effectifs ;
\item déterminer les fréquences si nécessaire ;
\item choisir une représentation adaptée ;
\item calculer une moyenne si elle répond à la question ;
\item interpréter les résultats ;
\item vérifier échelles, unités et cohérence.
\end{enumerate}
\end{methode}

\section{23. Erreurs fréquentes}

\begin{itemize}
\item Confondre effectif et fréquence.
\item Diviser par le mauvais effectif total.
\item Oublier de convertir une fréquence en pourcentage.
\item Construire un graphique sans titre ni échelle.
\item Utiliser un diagramme circulaire pour une évolution temporelle.
\item Calculer une moyenne en divisant par une mauvaise quantité de valeurs.
\item Interpréter une moyenne comme une valeur obligatoirement observée.
\item Lire une différence visuelle sans vérifier l'échelle.
\end{itemize}

\section{À retenir}

\begin{aretenir}
La fréquence vaut :
\[
\boxed{
f=\dfrac{\text{effectif}}{\text{effectif total}}
}.
\]

Elle peut être écrite sous forme fractionnaire, décimale ou en pourcentage.

Une série peut être représentée par un tableau, un diagramme ou un graphique.

La représentation doit être choisie en fonction de ce que l'on veut mettre en évidence.

La moyenne simple vaut :
\[
\boxed{
\overline{x}
=
\dfrac{\text{somme des valeurs}}{\text{nombre de valeurs}}
}.
\]

Toute représentation statistique doit être lue avec son contexte, ses unités et son échelle.
\end{aretenir}
